\documentclass[12pt]{article}

\usepackage[a4paper,top=2cm,bottom=2cm,left=2cm,right=2cm,marginparwidth=1.5cm]{geometry}
\usepackage{graphicx}
\usepackage{amssymb}
\usepackage{epstopdf}
\usepackage[german]{babel}
\usepackage[utf8]{inputenc}
\usepackage{pdfpages}
\usepackage{fancyhdr}
\DeclareGraphicsRule{.tif}{png}{.png}{`convert #1 `dirname #1`/`basename #1 .tif`.png}

\title{Bericht}
\author{Alisa Pesteritz, Saliha Altan, Wiktoria Cholewa, \\ 
	Dominik Falk, Andreas Kasarinow, Johanna Kleidorfer, \\
	Adrika Narasimhan, Rosina Röckseisen}

\begin{document}

\newpage
\maketitle

\newpage
\pagestyle{fancy}
\fancyhead{}
\fancyhead[L]{Alle}
\section{Einleitung}
\label{einleitung}

\section{Arbeitseinteilung (oder sowas)}
\label{arbeitseinteilung}
% TODO: wer hat was gemacht und warum?
% Datenfluss (von Johanna) einfügen? (oder zumindest beschreiben?)

\section{Grundlagen}
\label{grundlagen}

\subsection{Biologische Grundlagen}
\label{bio-grundlagen}

\subsubsection{Genexpression und -regulation}

Die im Folgenden analysierten Genexpressionsdaten beschreiben die Aktivität eines Gens und damit wann und in welcher Intensität ein Gen abgelesen wird. Dies entscheidet über den Phänotyp einer Zelle, denn auch wenn die Zellen eines Organismus dieselbe genetische Information in Form von DNA enthalten, erfüllen sie doch sehr unterschiedliche Funktionen. Die Genexpression beschreibt dabei das Zentrale Dogma der Molekularbiologie, von der Transkription wird die in der DNA codiert Information in RNA und anschließend mittels der Translation in Proteine überführt \cite{LeeYoung2013}.

Dieser Prozess der Genexpression unterliegt einer zellinternen Genregulation. Dabei binden Transkriptionsfaktoren an regulatorische DNA-Abschnitte wie Promotoren und Enhancer und rekrutieren drüber die RNA-Polymerase, welche das Ablesen des jeweiligen Genabschnitts startet. Je nach gebundenem Transkriptionsfaktor wird die Transkription gehemmt oder aktiviert, es kommt also entweder zum Ablesen oder nicht. Transkriptionsfaktoren sind dabei selbst Proteine, die als Genprodukt anderer Gene entstehen. Je nach Zelltypen, Entwicklungsstadien und äußeren Signalen sind verschiedene Transkiptionsfaktoren vorhanden und aktiv. Diese kaskadierenden Genregulationsnetzwerke erklären warum trotz identischer DNA in jeder Zelle andere Gene ausgelesen werden \cite{LeeYoung2013}.

Die Feinjustierung der Genexpression hat mehrere nachgelagerte Ebenen. Post-transkriptionelle wird beispielsweise durch microRNAs oder alternatives Spleißen reguliert, welche und wie viel mRNA für die Translation zur Verfügung steht. Translational wird die Proteinbiosyntheserate kontrolliert, welche bestimmt, wie viele funktionale Proteine aus einer gegebenen mRNA-Menge tatsächlich entstehen. In der postranslationalen Ebene findet die Proteinmodifikation (z.\,B. Phosphorylierung, Ubiquitinierung) statt. Eine Abnormalität in einer dieser Kontrollebenen kann zur Fehlregulation der Genexpression und in Folge dessen zu Krankheitsbildern führen \cite{LeeYoung2013}.

Die Datenerhebung der möglicherweise dafür verantwortlichen RNA-Mengen erfolgt nach heutigem Stand der Technik meist über das High-Throughput-Verfahren des Next-Generation-Sequencing (RNA-Seq). In Ausnahmefällen werden noch ältere Microarray-Technologien verwendet, diese haben allerdings eine geringere Sensitivität und einen eingeschränkteren dynamischen Messbereich als RNA-Seq \cite{Wang2009}. Der bioinformatische Workflow führt dabei von der RNA-Isolation über das Alignment an ein Referenzgenom bis hin zur Generierung einer rohen Count-Matrix \cite{Kukurba2015}. Diese Matrix bildet das Ergebnis der Sequenzierung ab, sie enthält für jedes Gen (Zeilen) und jede untersuchte Probe oder jeden Patienten (Spalten) die Anzahl der Reads. Eine solches tabellarische Format erwartet auch der CSV-Upload (siehe Kap.~\ref{gui}).

Bei diesen Rohdaten kommt es jedoch unweigerlich zu systematischen Verzerrungen, bedingt durch variierende Genlängen und stark schwankende Sequenziertiefen zwischen den einzelnen Probenläufen. Die Sequenziertiefe bezeichnet dabei die Gesamtzahl der in einem Lauf erzeugten Reads für eine Probe, welche stark variieren können. Ein direkter Vergleich zweier Proben unterschiedlicher Sequenziertiefen würde also fälschlicherweise einen Unterschied in der Genaktivität suggerieren, obwohl dieser lediglich an der Datenmenge liegt. Für eine valide biologische Vergleichbarkeit müssen die Datensätze daher zwingend mathematisch normalisiert werden, bevor sie miteinander verglichen werden können (siehe Kap.~\ref{normalisierung}) \cite{Oshlack2010}.

Die Analyse der Genexpressionsdaten ist vorrangig für die Onkologie von großem Interesse. Dabei wird die Genaktivität quantifiziert und analysiert, wofür das Transkriptom, also die Gesamtheit aller zu einem bestimmten Zeitpunkt in einer Zelle transkribierten RNA-Moleküle, betrachtet werden muss \cite{Conesa2016}. Denn nicht nur das Genom, also die strukturelle DNA und deren Mutationen, ist entscheidend für die Entstehung und das Fortschreiten vieler Krankheitsbilder, sondern maßgeblich auch die zelluläre Aktivität. So lassen sich beispielsweise pathophysiologische Prozesse wie unkontrolliertes Zellwachstum oder die Metastasierung bei strukturell intakten Genen beobachten, wenn diese fehlerhaft reguliert (über- oder unterexprimiert) sind \cite{Vogelstein2013}.

\subsubsection{Pathways und die KEGG-Datenbank}

Da viele biologische Prozesse jedoch nicht monogen ablaufen, reicht die Betrachtung einzelner Gene meist nicht für die biologische Interpretierbarkeit. Deshalb müssen funktionell zusammenhängende Gene gemeinsam betrachtet werden, um Muster und Zusammenhänge erkennen zu können. Dementsprechend ist es wichtig die Gene zusammen darzustellen, welche an denselben biologischen Signalwegen beteiligt sind, dies kann durch sogenannte Pathways realisiert werden. Dabei handelt es sich um kuratierte Netzwerke von Genen bzw. Genprodukten, die gemeinsam einen gewissen Signalweg abbilden, wie beispielsweise den einer Signaltransduktion. Die Pathways reduzieren für die Analyse der Daten die Komplexität und erhöhen dabei gleichzeitig die Interpretierbarkeit der Ergebnisse, da sie den funktionalen Kontext liefern \cite{Khatri2012}.

Eine der wichtigsten Referenzdatenbanken, welche hunderte dieser kuratierten Pathways für unterschiedlichste Organismen bereithält, ist die Kyoto Encyclopedia of Genes and Genomes (KEGG) \cite{Kanehisa2023}. KEGG bildet Signalwege, Stoffwechselwege und krankheitsassoziierte Prozesse als standardisierte grafische annotierte Netzwerke ab und beinhaltet pro Pathway die zusammengehörigen Gene. Damit können die Genexpressionsdaten direkt auf die bekannten biologischen Zusammenhänge abgebildet werden. 

Das in diesem Projekt entwickelte Dashboard erlaubt es den Nutzern verschiedenste Pathways der KEGG Datenbank auszuwählen. Damit können die interessanten Aktivitätsmuster gezielt und ohne Umwege über die vollständige Genexpressionsmatrix untersucht werden. Genaueres zur technischen Implementierung ist in Kapitel~\ref{pathways} zu finden.

\subsection{Medizinische Grundlagen}
\label{medizin-grundlagen}

\subsubsection{Klinische Relevanz der Genexpressionanalyse}

Das entwickelte Tumorboard Dashboard findet vorrangig im klinischen Kontext, bei der Entwicklung personalisierter Therapiepläne, Anwendung. Ärzte sind hierbei besonders an der Genexpression und deren Regulation interessiert, da herkömmliche histopathologische, also rein visuelle Untersuchungen des Tumorgewebes die molekulare Heterogenität von Tumoren meist nicht vollständig abbilden können \cite{Vogelstein2013}. Denn zwei histologisch identische Tumore können auf molekularer Ebene vollkommen unterschiedliche Aktivitätsprofile aufweisen und entsprechend unterschiedlich auf verschiedene Therapien ansprechen. Diese Informationsebene kommt im klinischen Bereich besonders für die Diagnosen, die Prognoseabschätzung und die Vorhersage der Erfolgschancen bei Therapieansätzen zum Einsatz \cite{NarrandesXu2018}. Beispiele aus der heutigen Zeit für den Einsatz der Genexpressionsdaten sind die Bestimmung des Ursprungsgewebes bei Tumoren unbekannter Herkunft, die Therapieentscheidung bei Brustkrebs sowie die Einschätzung des Rezidivrisikos, also der Wahrscheinlichkeit, dass die Erkrankung nach erfolgreicher Genesung erneut auftritt \cite{NarrandesXu2018}.

\subsubsection{Das Molekulare Tumorboard}

Das Molekulare Tumorboard (MTB) ist eine interdisziplinäre Expertenkonferenz, in der für besondere Krankheitsbilder oder spezielle Verläufe individualisierte Therapiepläne entwickelt werden sollen. Hier dienen die genomischen und transkriptomischen Profile der Entscheidungsfindung \cite{Tamborero2020}.

\subsubsection{Das Dashboard als Clinical Decision Support System}

Das Dashboard fungiert hierbei als interaktives Clinical Decision Support System (CDSS) \cite{Tamborero2020}. Es bietet mit der Heatmap und den Dendrogrammen eine gemeinsame, leicht verständliche, visuell dargestellte Diskussionsgrundlage für die beteiligten Experten. Es erlaubt einem interdisziplinären Kollektiv aus Onkologen, Pathologen, Chirurgen und Bioinformatikern die medizinischen Analysen ohne Programmierkenntnisse in Echtzeit über eine grafische Benutzeroberfläche zu steuern und anschließend gemeinsam darüber zu beraten. 

Das Ziel des interdisziplinären Gremiums besteht dementsprechend darin, die genomischen und transkriptomischen Profile der Patienten mithilfe des CDSS gemeinschaftlich zu evaluieren und mit den daraus generierten Informationen eine individuelle Therapieempfehlung auszuarbeiten. Denn wichtige und weitreichende Therapieentscheidungen können nicht von einem einzelnen Arzt isoliert verantwortet werden \cite{Mosele2020}. Studien belegen, dass diese Art der computergestützten Therapiewahl die Behandlungsergebnisse maßgeblich verbessern kann \cite{Horak2021}. Durch die Anpassungsmöglichkeiten, wie die Auswahl spezifischer biologischer Pathways, die Modifikation der Normalisierungsverfahren, die Wahl des Cluster-Algorithmus (z.\,B. Single, Average und Complete Linkage) und der farblichen Darstellung kann das Dashboard zum zentralen Visualisierungsmedium in der Konferenz werden. Das ermöglicht den Ärzten, Hypothesen direkt während der Sitzung gemeinsam zu prüfen, molekulare Muster im Kollektiv zu interpretieren und so die bestmögliche personalisierte Therapieempfehlung für die Patienten auszuarbeiten \cite{PereraBel2018}.

% ----------------------------------------------------------------------------------------------
\newpage
\pagestyle{fancy}
\fancyhead{}
\fancyhead[L]{Alisa}
\section{Pathways und Dateneinlesen}
\label{pathways}

% ----------------------------------------------------------------------------------------------
\newpage
\pagestyle{fancy}
\fancyhead{}
\fancyhead[L]{Rosina \& Wiki}
\section{Normalisierung}
\label{normalisierung}


\section{Clusteranalyse}
\label{clusteranalyse}
% TODO: theoretische Erklärung des Algorithmus

\subsection{Single Linkage}
\label{single-linkage}

\subsection{Average Linkage}
\label{average-linkage}

\subsection{Complete Linkage}
\label{complete-linkage}

% ----------------------------------------------------------------------------------------------
\newpage
\pagestyle{fancy}
\fancyhead{}
\fancyhead[L]{Saliha \& Domi}
\section{Visualisierung}
\label{visualisierung}

\subsection{Heatmap}
\label{heatmap}

\subsection{Dendrogramm}
\label{dendrogramm}

% ----------------------------------------------------------------------------------------------
\newpage
\pagestyle{fancy}
\fancyhead{}
\fancyhead[L]{Adrika \& Andreas}
\section{GUI}
\label{gui}

\subsection{Serverlogik}
\label{serverlogik}

% etc...

\newpage
\pagestyle{fancy}
\fancyhead{}
\fancyhead[L]{Johanna}
\subsection{Farben}


\end{document}
